\section{Problem Formulation and Method}
\label{sec:method}

\subsection{Patient-sequential adaptation}

Let $\mathcal{D}_0$ denote the source-patient (DS1) training data and
$f_{\theta_0}$ the source model. For an unseen target patient $p$, an early
chronological segment $\mathcal{D}^{\mathrm{adapt}}_p$ supplies labeled beats
for supervised adaptation, while a later segment $\mathcal{D}^{\mathrm{test}}_p$
is held out for evaluation: no test beat and no test label takes part in
adaptation or model selection. (The released preprocessing filters each record
as a whole, so the test \emph{signal} is not strictly unseen near the boundary;
Section~\ref{sec:protocol} bounds that dependency and Section~\ref{sec:splitfirst}
reruns the experiment without it.) Each patient begins from the \emph{same}
source checkpoint $\theta_0$; there is no cross-patient state carry-over. A
frozen replay memory $\mathcal{R}_0$ of source-domain exemplars is mixed into
adaptation to limit forgetting.

For a metric $\mathcal{M}$, the per-patient target gain and the per-patient
\emph{source-domain change} are
\begin{equation}
G_p =
\mathcal{M}\!\left(f_{\theta_p},\mathcal{D}^{\mathrm{test}}_p\right)
-\mathcal{M}\!\left(f_{\theta_0},\mathcal{D}^{\mathrm{test}}_p\right),
\qquad
\Delta^{\mathrm{src}}_p =
\mathcal{M}\!\left(f_{\theta_p},\mathcal{D}^{\mathrm{src}}_{\mathrm{test}}\right)
-\mathcal{M}\!\left(f_{\theta_0},\mathcal{D}^{\mathrm{src}}_{\mathrm{test}}\right).
\label{eq:gain_and_srcchange}
\end{equation}
Both are \emph{signed changes measured in the same direction}: positive means the
adapted model is better than the frozen source on that domain, negative means
worse. We therefore report $\Delta^{\mathrm{src}}_p$ throughout and call it the
\textbf{source-domain \macrof{} change}, not ``forgetting''. Every source-domain
number in this paper is negative, meaning source performance fell by that amount;
a value of $-0.027$ is a smaller loss than $-0.136$. We avoid the word forgetting
for the quantity itself because the two common conventions differ in sign and
mixing them is a standard source of confusion.
The primary unit of analysis is the patient: $\mathcal{M}$ is per-patient
\macrof{} over the AAMI classes present in that patient's test segment.
Fig.~\ref{fig:protocol} summarizes the adaptation protocol end to end.

\begin{figure}[H]
\centering
\includegraphics[width=\textwidth]{fig_method_protocol.pdf}
\caption{Method and experimental protocol. A DS1-trained tiny ECG transformer
is adapted per patient: source replay is mixed with the early (chronologically
first) labeled patient segment, while the later segment is held out for
testing: it contributes no beat, label, normalization statistic, early-stopping
signal, or model-selection information. The bounded whole-record filter
dependency near the chronological boundary is quantified and removed in the
split-first sensitivity analysis of Section~\ref{sec:splitfirst}. Adaptation may touch the classifier head, a post-pool residual
adapter, or an encoder-attention LoRA update, each drawing from the same
persistent-state budget. The budget covers persistent adaptation state; peak
temporary training memory and hardware runtime are not evaluated.}
\label{fig:protocol}
\end{figure}

\subsection{Analytical adaptation-state model}

\paragraph{What ``persistent state'' means here.} We use \emph{persistent
adaptation state} to mean \textbf{a statically reserved writable region that the
adaptation procedure needs across update steps}. It does \emph{not} imply that
every included quantity is nonvolatile or must survive power loss, and it is
\emph{not} total SRAM. The budget covers trainable weights, optimizer moments,
replay values, replay labels, and buffer metadata, all of which must persist
between updates. It also charges gradients, which are strictly update-time
working state rather than retained state; we include them because an
implementation that materializes a full gradient vector must reserve that region,
and excluding them would understate what a device has to set aside. It excludes
forward and backward activations, stack and workspace, and allocator overhead,
which are transient peak-SRAM quantities and are not measured anywhere in this
paper (Table~\ref{tab:memcats}).


A complete persistent-state accounting includes
\begin{equation}
B_{\mathrm{persistent}} =
B_{\mathrm{weights}}+B_{\mathrm{gradients}}+B_{\mathrm{optimizer}}
+B_{\mathrm{replay}}+B_{\mathrm{labels}}+B_{\mathrm{metadata}}+B_{\mathrm{padding}}.
\label{eq:completebudget}
\end{equation}
The experiments use FP32 trainable weights, FP32 gradients, and two FP32 Adam
moments, giving a per-parameter trainable cost of
\begin{equation}
b_{\mathrm{train}}=\underbrace{4}_{\text{weight}}+\underbrace{4}_{\text{grad}}
+\underbrace{8}_{\text{2 moments}}=16\ \mathrm{bytes/parameter}.
\end{equation}
A stored replay item costs more than the tensor it carries, and conflating the
two is an easy way to understate a budget. We therefore separate the
\emph{value payload} from the \emph{serialized record}:
\begin{equation}
b_l^{\mathrm{record}} =
b_l^{\mathrm{payload}} + b_{\mathrm{label}} + b_{\mathrm{flags}}
+ b_{\mathrm{pad}},
\label{eq:record}
\end{equation}
where $b_l^{\mathrm{payload}}$ is the INT8 tensor at replay location $l$ and the
remaining terms are the per-entry label, the validity and class-id flags, and any
alignment padding. Concretely
$b^{\mathrm{payload}} = 18\bytes \to b^{\mathrm{record}} = 21\bytes$ for
post-pool and $200\bytes \to 203\bytes$ for raw, at 1-byte alignment. For $P$
trainable parameters and $N$ replay items the reserved adaptation state is
\begin{equation}
B(P,N,l)=16P+N\,b_l^{\mathrm{record}}+B_{\mathrm{bufmeta}},
\label{eq:completebudget2}
\end{equation}
with $B_{\mathrm{bufmeta}}$ the fixed per-buffer overhead (write
index, count, reservoir state, class counters, and quantization parameters:
$49\bytes$ post-pool, $54\bytes$ raw). The constrained replay-count generator that
fixes a fair budget match is therefore
\begin{equation}
N_{\max}(P,l)=
\left\lfloor
\frac{B_{\mathrm{budget}}-16P-B_{\mathrm{bufmeta}}}
     {b_l^{\mathrm{record}}}
\right\rfloor .
\label{eq:nmax}
\end{equation}
Every byte total in this paper uses $b_l^{\mathrm{record}}$, never
$b_l^{\mathrm{payload}}$; where the text quotes 18 or 200 bytes it is naming the
payload explicitly.
Every configuration in this paper is generated from Eq.~\eqref{eq:nmax} so that
$B(P,N_{\max},l)\le 16{,}384\bytes$ holds by construction; the packed byte total
for each arm is asserted per adaptation cell. Table~\ref{tab:arm_definitions}
lists the resulting trainable counts, replay counts, and byte totals.

\paragraph{One-beat decision latency from the post-RR feature.} The model
consumes both the preceding and the \emph{following} RR interval. The following
interval is only known once the next R peak has been detected, so a classification
cannot be emitted at the current R peak: \textbf{the task evaluated here is
delayed heartbeat classification, one beat behind the signal}, not zero-latency
classification. This is a property of the source architecture we adopt rather
than of the adaptation method, and it applies equally to every arm including the
frozen baseline, so it does not affect any comparison in this paper. It does bound
the deployment claim, and comparing pre-RR-only, predicted-post-RR, and delayed
variants is left as future work.

\subsection{Replay-point cliff}

The base model accepts a 198-sample ECG window and two RR-derived interval
values. A convolutional embedding produces 66 tokens of width 16; the
feed-forward hidden dimension is 128; mean pooling then contracts the sequence,
and the two RR values are fused post-pool. Table~\ref{tab:replay_bytes} and
Fig.~\ref{fig:cliff} give the per-item INT8 payload at each candidate replay tap.
Only the post-pool (and post-pool$+$RR) representation is smaller than the raw
input; convolution tokens are $5.3\times$ and the FFN hidden tensor $42.7\times$
larger. The cliff is deterministic for this architecture and fixes \emph{where}
replay may be stored, but not \emph{where it should be}.

\begin{table}[H]
\centering
\small
\caption{Replay-point sizes and the \emph{serialized} INT8 cost of one stored
exemplar at each candidate tap. ``Values'' is the tensor itself; ``$+$RR'' adds
the two RR-interval features the model also consumes, without which an exemplar
cannot be replayed; ``$+$meta'' adds the per-entry label, valid flag, and class
id. The stride is the per-entry cost under 1-byte alignment, which is exact here
because every field is a byte. Only the post-pool tap is storage-compressive.}
\label{tab:replay_bytes}
\setlength{\tabcolsep}{5pt}
\begin{tabular}{llrrrrl}
\toprule
Replay point & Shape & Values & $+$RR & $+$meta & $\times$ raw & Verdict \\
\midrule
Raw ECG window          & $198\times1$  & 198   & 200   & 203   & $1.00$  & baseline \\
After conv embedding    & $66\times16$  & 1{,}056 & 1{,}058 & 1{,}061 & $5.33$  & larger (disqualified) \\
Encoder output          & $66\times16$  & 1{,}056 & 1{,}058 & 1{,}061 & $5.33$  & larger (disqualified) \\
FFN hidden              & $66\times128$ & 8{,}448 & 8{,}450 & 8{,}453 & $42.67$ & larger (disqualified) \\
Post-pool (mean)        & $16$          & 16    & 18    & 21    & $0.08$  & \textbf{smaller (viable)} \\
Post-pool $+$ RR concat & $18$          & 18    & 18    & 21    & $0.09$  & \textbf{smaller (viable)} \\
\bottomrule
\end{tabular}

\vspace{0.4em}
{\small Buffer-level fixed metadata is charged once, not per entry: a 4\,B write
index, a 4\,B count, 16\,B of reservoir state, $5\times4$\,B class counters, and a
4\,B scale plus 1\,B zero point per quantized tensor. That is 49\,B for a
post-pool buffer (one quantized tensor) and 54\,B for a raw buffer (two: the beat
and the RR pair). Under 4-byte alignment the strides become 24\,B and 204\,B;
we report the 1-byte figures and account for padding explicitly.}
\end{table}


\begin{figure}[H]
\centering
\includegraphics[width=0.82\textwidth]{fig_replay_cliff.pdf}
\caption{Replay-point cliff (log scale). Token and feed-forward representations
expand relative to the raw ECG window; only post-pool replay is
storage-compressive. Storing post-pool exemplars therefore maximizes replay
\emph{count}, which motivates---but does not justify---maximum-volume post-pool
replay.}
\label{fig:cliff}
\end{figure}

\subsection{Adaptation arms}

Six arms span the replay-volume / plasticity trade-off at 16\,KiB:
\begin{itemize}[leftmargin=*]
\item \textbf{A0} --- frozen source model (reference; 0 trainable, 0 replay);
\item \textbf{A1} --- 95-parameter linear head, maximum post-pool replay;
\item \textbf{A2} --- rank-1 post-pool residual adapter plus head, post-pool replay;
\item \textbf{A3} --- rank-2 post-pool residual adapter plus head, post-pool replay;
\item \textbf{A4} --- rank-1 encoder-attention LoRA plus head, raw replay;
\item \textbf{A5} --- rank-2 encoder-attention LoRA plus head, raw replay.
\end{itemize}
For fused representation $z\in\mathbb{R}^{18}$, the post-pool adapter applies
$z'=z+BAz$ with $A\in\mathbb{R}^{r\times18}$, $B\in\mathbb{R}^{18\times r}$. For
an encoder projection $W$, LoRA applies $W'=W+BA$ with
$A\in\mathbb{R}^{r\times d_{\mathrm{in}}}$, $B\in\mathbb{R}^{d_{\mathrm{out}}\times r}$.
A1--A3 spend the budget on \emph{volume}: 650--705 post-pool exemplars at
18\,B of payload each (21\,B serialized). A4--A5 spend it on \emph{depth}: only
52--62 raw exemplars at 200\,B payload (203\,B serialized) each, but expose the encoder to low-rank updates. For context,
full-model Adam-style trainable state would need $6643\times16=106{,}288$ bytes,
far beyond the ceiling, so full fine-tuning is infeasible under the constraint.
Fig.~\ref{fig:budget} shows how each arm partitions the same 16\,KiB.

\begin{table}[H]
\centering
\small
\caption{Complete arm definitions. ``Trainable tensors'' names exactly what
receives gradients; every other weight is frozen. Persistent bytes cover
trainable weights, gradients, two Adam moments per parameter, INT8 replay values,
per-entry label/valid/class-id metadata, buffer-level metadata, and alignment
padding. Replay counts are the maximum that fit, generated by the constrained
configuration generator. The right-hand block repeats the generation against an
effective ceiling of $15{,}360$\,B, i.e.\ a 1\,KiB firmware implementation
reserve; every arm was re-run at those reduced counts
(Section~\ref{sec:reserve}).}
\label{tab:arm_definitions}
\setlength{\tabcolsep}{4pt}
\begin{tabular}{llp{3.5cm}rrrrr}
\toprule
& & & & \multicolumn{2}{c}{16{,}384\,B ceiling} & \multicolumn{2}{c}{$+$1\,KiB reserve} \\
\cmidrule(lr){5-6}\cmidrule(lr){7-8}
Arm & Replay & Trainable tensors & Params & Items & Bytes & Items & Bytes \\
\midrule
A0 & ---       & none (frozen)                                   & 0   & 0   & 0        & 0   & 0 \\
A1 & post-pool & head $W,b$                                      & 95  & 705 & 16{,}374 & 656 & 15{,}345 \\
A2 & post-pool & post-pool adapter $B_1A_1$ $+$ head             & 131 & 678 & 16{,}383 & 629 & 15{,}354 \\
A3 & post-pool & post-pool adapter $B_2A_2$ $+$ head             & 167 & 650 & 16{,}371 & 601 & 15{,}342 \\
A4 & raw       & encoder LoRA $r{=}1$ on $W_{q,k,v,o}$ $+$ head  & 223 & 62  & 16{,}208 & 57  & 15{,}193 \\
A5 & raw       & encoder LoRA $r{=}2$ on $W_{q,k,v,o}$ $+$ head  & 351 & 52  & 16{,}226 & 47  & 15{,}211 \\
\midrule
R1 & ---       & encoder LoRA $r{=}1$ $+$ head, L2 source anchor & 223 & 0   & 3{,}568  & 0   & 3{,}568 \\
R2 & raw       & encoder LoRA $r{=}1$ $+$ head, L2 source anchor & 223 & 62  & 16{,}208 & 57  & 15{,}193 \\
\bottomrule
\end{tabular}

\vspace{0.4em}
{\small All arms use Adam ($b_{\mathrm{train}} = 16$\,B/parameter: 4 weight,
4 gradient, 8 for two moments), FP32 weights, and INT8 replay. Replay and target
beats are concatenated rather than mixed at a fixed per-batch ratio, so the
effective replay fraction is $N_{\mathrm{replay}}/(N_{\mathrm{replay}} +
N_{\mathrm{target}})$ with $N_{\mathrm{target}} = 500$. The R-arm source anchor
adds \emph{no} persistent bytes: the anchor is the frozen base weight, which is
already stored for inference.}
\end{table}


\begin{figure}[H]
\centering
\includegraphics[width=0.80\textwidth]{fig_budget_composition.pdf}
\caption{Persistent-state composition of each arm at 16\,KiB. Volume arms
(A1--A3) allocate almost the entire budget to replay values; depth arms
(A4--A5) shift bytes into trainable weights, gradients, and optimizer moments,
leaving room for only tens of raw exemplars.}
\label{fig:budget}
\end{figure}
